\documentclass[12pt,a4paper]{article}

\usepackage[T2A]{fontenc}
\usepackage[utf8]{inputenc}
\usepackage[russian]{babel}
\usepackage{amsmath,amssymb}
\usepackage{booktabs}
\usepackage{longtable}
\usepackage{array}
\usepackage{geometry}
\usepackage{enumitem}

\geometry{margin=2.2cm}

\title{Исследовательские гипотезы, проверяемость, валидность и подтверждаемость моделей ИИС}
\author{Исследовательский проект по интегральному индексу состояния}
\date{\today}

\begin{document}
\maketitle

\section{Назначение документа}
Данный документ формулирует исследовательские гипотезы, связанные с версиями 1--4 модели ИИС, и отдельно оценивает:
\begin{enumerate}[leftmargin=1.25cm]
    \item \textbf{проверяемость},
    \item \textbf{валидность},
    \item \textbf{подтверждаемость}
\end{enumerate}
каждой версии на основе уже выполненных прогонов по реальным открытым физиологическим датасетам.

\section{Исходные данные для проверки}
На текущем этапе проекта фактически использовались четыре набора:
\begin{itemize}[leftmargin=1.25cm]
    \item \textbf{WESAD};
    \item \textbf{CASE};
    \item \textbf{OpenNeuro ds002722};
    \item \textbf{OpenNeuro ds002724}.
\end{itemize}

Из них для реального сравнения математических версий ИИС наиболее важны именно ds002722 и ds002724, потому что:
\begin{itemize}[leftmargin=1.25cm]
    \item содержат EEG;
    \item содержат ECG;
    \item содержат self-report valence/arousal;
    \item позволяют проверять блоки $A$, $\Gamma$, $V$, $Q$.
\end{itemize}

WESAD и CASE важны как дополнительные контрольные наборы:
\begin{itemize}[leftmargin=1.25cm]
    \item WESAD хорошо задаёт сценарий baseline/stress/amusement;
    \item CASE полезен для анализа временной динамики;
    \item однако оба набора не содержат EEG и потому слабо подходят для сравнения полных формул версий 1--4.
\end{itemize}

\section{Исследовательские гипотезы}
\subsection{Гипотеза H1}
\textbf{Разные версии ИИС ведут себя по-разному на одинаковых реальных данных.}

Смысл гипотезы: модели 1--4 должны давать не идентичные значения, а различимое поведение по разделению baseline/stress/disbalance.

\subsection{Гипотеза H2}
\textbf{Версии ИИС, использующие EEG и вегетативные признаки, способны разделять эмоциональные состояния на реальных открытых датасетах.}

Смысл гипотезы: если модель содержит физиологически содержательные блоки, то распределения IIS для baseline и stress должны различаться статистически и визуально.

\subsection{Гипотеза H3}
\textbf{Гормональный блок делает модель сильнее по разделению состояний, но одновременно снижает её методическую чистоту на открытых данных.}

Смысл гипотезы: версии 1--3 могут выигрывать за счёт proxy-гормонального блока, но их физиологическая интерпретация зависит от непрямых переменных.

\subsection{Гипотеза H4}
\textbf{Версия 4 без гормонов является более чистой эмоциональной моделью, хотя в текущем виде может уступать версиям 1--3 по силе разделения.}

Смысл гипотезы: отказ от труднодоступных биохимических переменных должен сделать модель честнее и удобнее для дальнейшей независимой проверки.

\subsection{Гипотеза H5}
\textbf{Компонент $Q$ является ключевым интегратором эмоционального состояния.}

Смысл гипотезы: если $Q$ действительно интегрально отражает качество состояния, то он должен быть связан с внешними показателями valence/arousal и оказывать заметный вклад в итоговый IIS.

\section{Что здесь понимается под проверяемостью, валидностью и подтверждаемостью}
\subsection{Проверяемость}
Под \textbf{проверяемостью} понимается принципиальная возможность вычислить модель на реальных данных и сравнить её прогнозы с внешними признаками состояния.

\subsection{Валидность}
Под \textbf{валидностью} понимается содержательная адекватность модели: отражает ли она реально измеряемые свойства состояния, а не произвольную комбинацию чисел.

\subsection{Подтверждаемость}
Под \textbf{подтверждаемостью} понимается наличие эмпирических результатов, которые поддерживают гипотезу, даже если они ещё не являются окончательным доказательством.

\section{Проверяемость версий}
\subsection{Версия 1}
\textbf{Статус: высокая проверяемость в proxy-режиме.}

Основания:
\begin{itemize}[leftmargin=1.25cm]
    \item компоненты $A$, $\Gamma$, $V$ и диагностический $Q$ считаются по реальным признакам;
    \item итоговый IIS рассчитывается на ds002722 и ds002724 без технических проблем;
    \item ограничение связано только с тем, что $H$ опирается на proxy-переменные.
\end{itemize}

\subsection{Версия 2}
\textbf{Статус: высокая проверяемость в proxy-режиме.}

Основания:
\begin{itemize}[leftmargin=1.25cm]
    \item вычислительно версия 2 столь же проверяема, как версия 1;
    \item все необходимые данные для $A$, $\Gamma$, $V$ и $Q$ реально доступны на ds002722 и ds002724.
\end{itemize}

\subsection{Версия 3}
\textbf{Статус: средняя или высокая проверяемость, но только в proxy-режиме.}

Основания:
\begin{itemize}[leftmargin=1.25cm]
    \item версия 3 формально вычисляется на реальных данных;
    \item однако блоки $K$, $DA$ и $C$ не имеют прямых измерений и требуют proxy-логики;
    \item поэтому модель проверяема как вычислительная, но не как прямой биохимический тест.
\end{itemize}

\subsection{Версия 4}
\textbf{Статус: высокая проверяемость как эмоциональной модели.}

Основания:
\begin{itemize}[leftmargin=1.25cm]
    \item блоки $A$, $\Gamma$, $V$ вычисляются по реальным EEG/ECG/HRV-признакам;
    \item гормональные блоки удалены;
    \item итоговый IIS проверяется на реальных данных;
    \item ограничение состоит в том, что в proxy-режиме компонент $Q$ частично якорится на self-report.
\end{itemize}

\section{Валидность версий}
\subsection{Версии 1 и 2}
\textbf{Статус: рабочая операциональная валидность, но неполная биологическая валидность.}

Это означает следующее:
\begin{itemize}[leftmargin=1.25cm]
    \item модели хорошо разделяют состояния;
    \item EEG- и HRV-блоки валидны как физиологические признаки;
    \item но гормональный блок на открытых данных остаётся непрямым, поэтому биологическая интерпретация $H$ ограничена.
\end{itemize}

\subsection{Версия 3}
\textbf{Статус: средняя валидность.}

Сильная сторона версии 3:
\begin{itemize}[leftmargin=1.25cm]
    \item идея модулирующего коэффициента $K$ физиологически правдоподобна.
\end{itemize}

Слабая сторона версии 3:
\begin{itemize}[leftmargin=1.25cm]
    \item биологический смысл $K$ пока не подтверждён прямыми измерениями;
    \item на открытых данных $K$ остаётся латентным расчётным коэффициентом, а не биомаркером.
\end{itemize}

\subsection{Версия 4}
\textbf{Статус: самая чистая валидность для эмоционального состояния.}

Это связано с тем, что:
\begin{itemize}[leftmargin=1.25cm]
    \item версия 4 опирается только на реально доступные EEG/ECG/HRV-признаки;
    \item гормональные предположения из неё убраны;
    \item она лучше соответствует задаче оценки текущего эмоционального состояния.
\end{itemize}

Ограничение:
\begin{itemize}[leftmargin=1.25cm]
    \item в proxy-режиме компонент $Q$ частично использует self-report, поэтому на текущем этапе версия 4 ещё не является полностью автономной физиологической моделью.
\end{itemize}

\section{Подтверждаемость гипотез по результатам}
\subsection{Сводная таблица}
\begin{longtable}{>{\raggedright\arraybackslash}p{2.7cm}cccccc}
\toprule
Версия & Датасет & Effect size & Overlap & Corr$(IIS, arousal)$ & Corr$(Q, valence)$ & Надёжность \\
\midrule
\endfirsthead
\toprule
Версия & Датасет & Effect size & Overlap & Corr$(IIS, arousal)$ & Corr$(Q, valence)$ & Надёжность \\
\midrule
\endhead
V1 & ds002722 & $-2.410$ & $0.212$ & $-0.280$ & $-0.017$ & high \\
V2 & ds002722 & $-2.412$ & $0.224$ & $-0.281$ & $-0.017$ & high \\
V3 & ds002722 & $-1.853$ & $0.329$ & $-0.273$ & $-0.017$ & high \\
V4 & ds002722 & $-0.959$ & $0.522$ & $-0.199$ & $0.424$ & medium \\
\midrule
V1 & ds002724 & $-2.687$ & $0.174$ & $-0.180$ & $-0.091$ & high \\
V2 & ds002724 & $-2.693$ & $0.174$ & $-0.180$ & $-0.091$ & high \\
V3 & ds002724 & $-2.178$ & $0.273$ & $-0.185$ & $-0.091$ & high \\
V4 & ds002724 & $-1.446$ & $0.455$ & $-0.233$ & $0.493$ & medium \\
\bottomrule
\end{longtable}

\subsection{Подтверждаемость гипотезы H1}
\textbf{Гипотеза H1 подтверждается.}

Основание:
\begin{itemize}[leftmargin=1.25cm]
    \item версии 1--4 дают разные effect size, разные overlap и разные ранги полезности;
    \item это означает, что на одинаковых данных модели не схлопываются в одно и то же поведение.
\end{itemize}

\subsection{Подтверждаемость гипотезы H2}
\textbf{Гипотеза H2 подтверждается.}

Основание:
\begin{itemize}[leftmargin=1.25cm]
    \item на ds002722 и ds002724 версии 1--4 различают baseline и stress;
    \item для версий 1--3 эффект разделения очень сильный;
    \item для версии 4 эффект слабее, но всё ещё отчётливо наблюдается.
\end{itemize}

\subsection{Подтверждаемость гипотезы H3}
\textbf{Гипотеза H3 подтверждается частично.}

Основание:
\begin{itemize}[leftmargin=1.25cm]
    \item версии 1--3 действительно сильнее по разделению состояний, чем версия 4;
    \item однако это преимущество достигается ценой использования гормональных proxy-переменных;
    \item следовательно, усиление разделения сопровождается снижением биологической строгости на открытых данных.
\end{itemize}

\subsection{Подтверждаемость гипотезы H4}
\textbf{Гипотеза H4 подтверждается.}

Основание:
\begin{itemize}[leftmargin=1.25cm]
    \item версия 4 действительно является более чистой эмоциональной моделью;
    \item при этом она действительно уступает версиям 1--3 по силе разделения.
\end{itemize}

\subsection{Подтверждаемость гипотезы H5}
\textbf{Гипотеза H5 подтверждается частично.}

Основание:
\begin{itemize}[leftmargin=1.25cm]
    \item в версии 4 корреляция $Q$ с valence положительна и заметна: $0.424$ на ds002722 и $0.493$ на ds002724;
    \item значит, $Q$ действительно работает как интеграционный эмоциональный блок;
    \item однако часть этой связи обеспечивается proxy-якоризацией, поэтому полная независимая физиологическая подтверждаемость ещё не достигнута.
\end{itemize}

\section{Ограничения подтверждаемости}
\subsection{WESAD}
WESAD полезен для стрессового сценария, но не даёт EEG. Поэтому:
\begin{itemize}[leftmargin=1.25cm]
    \item версии 1--3 на нём почти схлопываются к HRV-блоку;
    \item полная проверка блоков $A$, $\Gamma$, $Q$ на WESAD невозможна.
\end{itemize}

\subsection{CASE}
CASE полезен для динамического анализа, но также не даёт EEG. Поэтому:
\begin{itemize}[leftmargin=1.25cm]
    \item он важен для траекторий и устойчивости;
    \item но недостаточен для строгого сравнения всех статических формул.
\end{itemize}

\subsection{Открытые ограничения}
Главные текущие ограничения проекта:
\begin{itemize}[leftmargin=1.25cm]
    \item нет прямых измерений гормонов;
    \item нет полностью автономной строгой версии 4 без self-report-якоря в $Q$;
    \item формальная динамическая математическая версия ИИС пока ещё не выделена в отдельную модель.
\end{itemize}

\section{Итоговые формулировки для проекта}
\subsection{Строгая формулировка}
На реальных открытых EEG+ECG-датасетах ds002722 и ds002724 все четыре версии ИИС оказались вычислимы и эмпирически сравнимы. Версии 1--3 обеспечивают более сильное разделение baseline и stress, однако опираются на proxy-гормональный блок. Версия 4, не использующая гормоны, является более методически чистой эмоциональной моделью и демонстрирует умеренное, но воспроизводимое разделение состояний.

\subsection{Формулировка про подтверждаемость}
Исследовательские гипотезы о различимости версий ИИС, о способности EEG+ECG-моделей разделять состояния и о преимуществе гормон-независимой версии для эмоциональной оценки получили эмпирическую поддержку. При этом гипотезы о прямом биохимическом смысле гормональных proxy-блоков и коэффициента $K$ пока нельзя считать окончательно доказанными.

\section{Вывод}
В текущем состоянии проекта ИИС уже представляет собой не только теоретическую схему, но и набор реально проверяемых моделей. Наиболее сильная с исследовательской точки зрения линия дальнейшей работы --- развитие версии 4 в сторону полностью автономной строгой эмоциональной модели без зависимости от self-report при сохранении проверяемости на реальных физиологических данных.

\end{document}
